
\section{Advanced Activities}

The following activities are intended for strong master's or PhD students.
They deepen the central ideas of this chapter: plant--control coupling,
system-level objectives, the performance loss caused by sequential design,
control authority, implementation realism, and the conditions under which
control co-design is worth the additional modeling and computational effort.

\subsection*{Problem 1: Exact Sequential versus Simultaneous Co-Design}

Consider the coupled quadratic design problem
\begin{equation}
J(p,c)
=
4p^2+3c^2+4pc-12p-10c+20,
\end{equation}
where $p$ is a plant-design variable and $c$ is a controller-design
variable. The design bounds are
\begin{equation}
0\leq p\leq3,
\qquad
0\leq c\leq3.
\end{equation}

\begin{enumerate}
    \item Compute the Hessian of $J$ and determine whether the problem is
    strictly convex.

    \item In a single-pass sequential workflow, set the nominal controller
    to
    \begin{equation}
    \hat{c}=0.
    \end{equation}
    First optimize $p$, then freeze $p$ and optimize $c$. Compute
    \begin{equation}
    p_{\mathrm{seq}},
    \qquad
    c_{\mathrm{seq}},
    \qquad
    J_{\mathrm{seq}}.
    \end{equation}

    \item Solve the simultaneous CCD problem
    \begin{equation}
    \min_{p,c}J(p,c)
    \end{equation}
    analytically.

    \item Compute the percentage improvement
    \begin{equation}
    \eta
    =
    \frac{
    J_{\mathrm{seq}}-J_{\mathrm{CCD}}
    }{
    J_{\mathrm{seq}}
    }
    \times100\%.
    \end{equation}

    \item Repeat the sequential calculation using
    \begin{equation}
    \hat{c}=0.5,\ 1,\ 2.
    \end{equation}
    Determine how strongly the final sequential design depends on the
    assumed nominal controller.

    \item Derive the value of $\hat{c}$ for which the single-pass sequential
    workflow reproduces the simultaneous optimum exactly.

    \item Explain why this special value does not make sequential design
    generally equivalent to CCD.
\end{enumerate}


\subsection*{Problem 2: Quantifying Plant--Control Coupling}

Consider the parameterized system-level objective
\begin{equation}
J(p,c;\gamma)
=
\frac{1}{2}a p^2
+
\frac{1}{2}b c^2
+
\gamma pc
-
r_p p
-
r_c c,
\end{equation}
where
\begin{equation}
a>0,
\qquad
b>0.
\end{equation}

The scalar $\gamma$ represents plant--control coupling strength.

\begin{enumerate}
    \item Derive the condition on $\gamma$ for $J$ to be strictly convex.

    \item Derive the optimal controller for a fixed plant:
    \begin{equation}
    c^*(p).
    \end{equation}

    \item Compute the controller sensitivity
    \begin{equation}
    \frac{dc^*}{dp}.
    \end{equation}

    \item Derive the optimal plant for a fixed controller:
    \begin{equation}
    p^*(c),
    \end{equation}
    and compute
    \begin{equation}
    \frac{dp^*}{dc}.
    \end{equation}

    \item Show that both sensitivities vanish when
    \begin{equation}
    \gamma=0.
    \end{equation}

    \item Solve for the simultaneous optimum
    \begin{equation}
    p^*(\gamma),
    \qquad
    c^*(\gamma).
    \end{equation}

    \item Use
    \begin{equation}
    a=4,
    \qquad
    b=3,
    \qquad
    r_p=6,
    \qquad
    r_c=5,
    \end{equation}
    and evaluate the optimum for
    \begin{equation}
    \gamma
    \in
    \{0,\ 0.5,\ 1,\ 2,\ 3\}.
    \end{equation}

    \item Plot
    \begin{equation}
    \left|
    \frac{dc^*}{dp}
    \right|
    \end{equation}
    and the sequential performance loss as functions of $\gamma$.

    \item Propose a dimensionless coupling index based on the Hessian
    cross-term and justify your choice.
\end{enumerate}


\subsection*{Problem 3: Control Authority and the Value of Integration}

Consider
\begin{equation}
m\ddot{x}+c\dot{x}+kx=u,
\end{equation}
with
\begin{equation}
m=1,
\qquad
x(0)=1,
\qquad
\dot{x}(0)=0.
\end{equation}

Use the PD controller
\begin{equation}
u=-K_px-K_d\dot{x},
\end{equation}
with
\begin{equation}
0\leq K_p\leq K_{p,\max},
\qquad
0\leq K_d\leq K_{d,\max}.
\end{equation}

The plant-design variables satisfy
\begin{equation}
0.1\leq k\leq6,
\qquad
0.05\leq c\leq3.
\end{equation}

Define
\begin{equation}
A_{\mathrm{cl}}
=
\begin{bmatrix}
0&1\\
-(k+K_p)&-(c+K_d)
\end{bmatrix}.
\end{equation}

For stable designs, let $P$ solve
\begin{equation}
A_{\mathrm{cl}}^TP
+
PA_{\mathrm{cl}}
+
Q
+
K^TRK
=
0,
\end{equation}
where
\begin{equation}
Q=
\begin{bmatrix}
1&0\\
0&0.05
\end{bmatrix},
\qquad
R=0.02,
\qquad
K=
\begin{bmatrix}
K_p&K_d
\end{bmatrix}.
\end{equation}

Define the total cost
\begin{equation}
J
=
\mathbf{x}_0^TP\mathbf{x}_0
+
0.02k^2
+
0.03c^2,
\qquad
\mathbf{x}_0=
\begin{bmatrix}
1\\0
\end{bmatrix}.
\end{equation}

\begin{enumerate}
    \item Derive the closed-loop stability conditions.

    \item Explain why
    \begin{equation}
    \mathbf{x}_0^TP\mathbf{x}_0
    \end{equation}
    equals the infinite-horizon dynamic performance cost.

    \item For
    \begin{equation}
    K_{p,\max}
    =
    K_{d,\max}
    \in
    \{0,\ 0.5,\ 1,\ 2,\ 4,\ 8\},
    \end{equation}
    solve:
    \begin{enumerate}
        \item the passive-first sequential problem;
        \item the simultaneous CCD problem.
    \end{enumerate}

    \item Plot the relative CCD advantage
    \begin{equation}
    \eta_{\mathrm{CCD}}
    =
    \frac{
    J_{\mathrm{seq}}-J_{\mathrm{CCD}}
    }{
    J_{\mathrm{seq}}
    }
    \end{equation}
    against allowable control authority.

    \item Determine whether the coupling becomes weak as the allowable
    controller gains approach zero.

    \item Determine whether unlimited control authority drives the optimal
    passive stiffness and damping to their lower bounds.

    \item Add the actuator-effort regularization
    \begin{equation}
    \beta(K_p^2+K_d^2)
    \end{equation}
    with
    \begin{equation}
    \beta\in\{0,\ 10^{-3},\ 10^{-2},\ 10^{-1}\},
    \end{equation}
    and repeat the study.

    \item Explain why strong control authority does not automatically imply
    that the most active design is the most practical design.
\end{enumerate}


\subsection*{Problem 4: System-Level Objective and Hidden Tradeoffs}

An active suspension study uses the weighted objective
\begin{equation}
J
=
w_aJ_a
+
w_tJ_t
+
w_uJ_u
+
w_mJ_m,
\end{equation}
where
\begin{align}
J_a&=\int_0^T a_s(t)^2dt,\\
J_t&=\int_0^T \Delta z_t(t)^2dt,\\
J_u&=\int_0^T F(t)^2dt,\\
J_m&=M(\mathbf{x}_p).
\end{align}

Two feasible designs produce the normalized metrics
\begin{equation}
\mathbf{J}^{(A)}
=
\begin{bmatrix}
0.60\\
1.10\\
1.45\\
0.75
\end{bmatrix},
\qquad
\mathbf{J}^{(B)}
=
\begin{bmatrix}
0.85\\
0.80\\
0.70\\
1.10
\end{bmatrix}.
\end{equation}

\begin{enumerate}
    \item Determine the region of nonnegative weight vectors
    \begin{equation}
    \mathbf{w}
    =
    \begin{bmatrix}
    w_a&w_t&w_u&w_m
    \end{bmatrix}^{T},
    \qquad
    \sum_iw_i=1,
    \end{equation}
    for which design $A$ is preferred.

    \item Determine the region for which design $B$ is preferred.

    \item Show that neither design dominates the other in the Pareto sense.

    \item Use the weight vectors
    \begin{align}
    \mathbf{w}^{(1)}
    &=
    \begin{bmatrix}
    0.55&0.20&0.15&0.10
    \end{bmatrix}^{T},
    \\
    \mathbf{w}^{(2)}
    &=
    \begin{bmatrix}
    0.25&0.25&0.25&0.25
    \end{bmatrix}^{T},
    \\
    \mathbf{w}^{(3)}
    &=
    \begin{bmatrix}
    0.20&0.25&0.15&0.40
    \end{bmatrix}^{T},
    \end{align}
    and determine which design is selected in each case.

    \item Suppose passenger comfort must satisfy the hard requirement
    \begin{equation}
    J_a\leq0.75.
    \end{equation}
    Determine which designs remain feasible.

    \item Suppose actuator effort must satisfy
    \begin{equation}
    J_u\leq1.
    \end{equation}
    Determine which designs remain feasible.

    \item Explain why moving an engineering requirement from the objective
    to a constraint can change the selected CCD solution even when the same
    physical metric is used.

    \item Construct a third feasible design $C$ that dominates either $A$
    or $B$ but not both.

    \item Explain why reporting only the aggregate weighted objective can
    hide an unacceptable engineering compromise.
\end{enumerate}


\subsection*{Problem 5: Ideal Actuator versus Implementable Actuator}

Consider the plant
\begin{equation}
\dot{x}=-x+f_a+d(t),
\qquad
x(0)=0,
\end{equation}
where
\begin{equation}
d(t)=\sin(4t).
\end{equation}

An idealized study assumes
\begin{equation}
f_a(t)=u_c(t).
\end{equation}

A more realistic actuator satisfies
\begin{equation}
\dot{f}_a
=
\omega_a
\left(
u_c-f_a
\right),
\end{equation}
with
\begin{equation}
|u_c(t)|\leq F_{\max}.
\end{equation}

The actuator design variables are
\begin{equation}
0.5\leq F_{\max}\leq5,
\qquad
2\leq\omega_a\leq50.
\end{equation}

Minimize
\begin{equation}
J
=
\int_0^{10}
\left(
x(t)^2
+
0.02u_c(t)^2
\right)dt
+
0.05F_{\max}^2
+
\frac{0.1}{\omega_a}.
\end{equation}

\begin{enumerate}
    \item For the ideal actuator, derive the control command that would
    cancel the disturbance exactly while maintaining
    \begin{equation}
    x(t)=0.
    \end{equation}

    \item Determine the minimum ideal actuator capacity required for exact
    disturbance cancellation.

    \item For the finite-bandwidth actuator, derive the steady-state
    sinusoidal amplitude and phase of $f_a(t)$ when
    \begin{equation}
    u_c(t)=-\sin(4t).
    \end{equation}

    \item Derive the resulting steady-state amplitude of $x(t)$.

    \item Evaluate the state-response amplitude for
    \begin{equation}
    \omega_a
    \in
    \{2,\ 4,\ 8,\ 16,\ 32,\ 50\}.
    \end{equation}

    \item Determine the smallest actuator bandwidth for which the state
    amplitude is within $5\%$ of the ideal-actuator prediction.

    \item Optimize $F_{\max}$ and $\omega_a$ for the realistic model using a
    suitable numerical method.

    \item Explain why a low objective obtained with the ideal actuator model
    is not sufficient evidence of a useful CCD result.

    \item Identify at least three other implementation effects that could
    invalidate an apparently strong CCD solution.
\end{enumerate}


\subsection*{Problem 6: When Is Control Co-Design Worth Using?}

A design team may choose either a sequential process or a CCD process.

For project $i$, define:
\begin{equation}
V_i
=
\text{expected lifecycle value of performance improvement},
\end{equation}
\begin{equation}
C_i
=
\text{additional modeling and computational cost of CCD},
\end{equation}
and
\begin{equation}
R_i
=
\text{expected cost of model or implementation failure}.
\end{equation}

The expected net value of CCD is
\begin{equation}
N_i
=
V_i-C_i-R_i.
\end{equation}

For a stylized dynamic system, assume
\begin{equation}
V
=
V_0
\left[
1-
\exp
\left(
-\alpha A\Gamma
\right)
\right],
\end{equation}
where:
\begin{equation}
A\geq0
\end{equation}
is normalized control authority and
\begin{equation}
\Gamma\geq0
\end{equation}
is normalized plant--control coupling strength.

Assume
\begin{equation}
C
=
C_0+C_1n_pn_c+C_2N_s,
\end{equation}
where $n_p$ and $n_c$ are the numbers of plant and controller variables and
$N_s$ is the number of uncertainty scenarios.

Assume
\begin{equation}
R
=
R_0
\exp
\left(
-\beta F
\right),
\end{equation}
where
\begin{equation}
0\leq F\leq1
\end{equation}
is a normalized model-fidelity and validation score.

Use
\begin{equation}
V_0=100,
\qquad
\alpha=0.8,
\qquad
C_0=5,
\end{equation}
\begin{equation}
C_1=0.15,
\qquad
C_2=0.5,
\qquad
R_0=40,
\qquad
\beta=3.
\end{equation}

\begin{enumerate}
    \item Derive the break-even condition
    \begin{equation}
    N=0.
    \end{equation}

    \item For
    \begin{equation}
    n_p=8,
    \qquad
    n_c=5,
    \qquad
    N_s=10,
    \end{equation}
    compute the minimum product
    \begin{equation}
    A\Gamma
    \end{equation}
    required for CCD to have positive expected net value when
    \begin{equation}
    F=0.8.
    \end{equation}

    \item Repeat for
    \begin{equation}
    F=0.3,\ 0.5,\ 1.
    \end{equation}

    \item Plot the break-even boundary in the $(A,\Gamma)$ plane.

    \item Determine how the break-even boundary changes when the number of
    uncertainty scenarios increases.

    \item Explain why high coupling and high control authority favor CCD.

    \item Explain why poor model fidelity can eliminate the expected value
    of CCD even when the numerical optimization predicts a large
    performance improvement.

    \item Use the model to compare the following projects:
    \begin{enumerate}
        \item a legacy machine with nearly fixed hardware and weak control
        authority;
        \item a flexible robot with strong sensing and actuation;
        \item a new floating wind platform with uncertain hydrodynamics;
        \item a regulated product with fixed plant architecture.
    \end{enumerate}

    \item Discuss the limitations of using a single scalar net-value model to
    decide whether CCD should be pursued.
\end{enumerate}
